\chapter{Balance Tests}

\section*{What Is This Test?}
Balance tests evaluate how the inner ear, eyes, brain, nerves, muscles, and joints work together to maintain posture and orientation. They help investigate dizziness, vertigo, unsteadiness, and falls.

\section*{Alternative Names}
\begin{itemize}
  \item Vestibular Testing
  \item Balance Function Testing
  \item Posturography
  \item Dizziness Evaluation
\end{itemize}

\section*{Principle}
The assessment may combine bedside eye and balance examinations with videonystagmography, caloric testing, rotary-chair testing, vestibular-evoked myogenic potentials, head-impulse testing, or computerized dynamic posturography. Each examines a different part of the balance system.

\section*{Why This Test Is Ordered}
Testing is used for recurrent vertigo, motion-provoked dizziness, unexplained imbalance, falls, suspected inner-ear disease, or symptoms after head injury. Hearing tests and neurologic assessment often provide important context.

\section*{Primary vs Secondary Test}
History and physical examination are primary. Specialized vestibular tests are secondary studies selected according to the suspected cause; no single balance test evaluates every part of the system.

\section*{How This Test Is Performed}
The patient may follow moving targets, wear eye-tracking goggles, move the head or body in controlled positions, stand on a force platform, or receive warm and cool air or water in the ear canal. Testing is performed by trained staff and may provoke brief dizziness or nausea.

\section*{What Preparation Is Needed}
Follow the testing center's instructions. Some medicines, alcohol, caffeine, and nicotine can affect results; do not stop prescribed medicines unless the clinician directs you to do so. Arrange help getting home if dizziness is expected.

\section*{Understanding Results}
Results may suggest a peripheral vestibular disorder, a central neurologic pattern, visual or sensory dependence, or no measurable deficit. Interpretation requires the symptoms, examination, hearing findings, and test protocol.

\section*{Follow-up Tests}
Follow-up may include audiometry, brain or inner-ear imaging, neurologic evaluation, orthostatic vital signs, physical therapy assessment, or targeted vestibular rehabilitation.

\section*{References}
\begin{enumerate}
  \item MedlinePlus. Balance Tests.
  \item American Academy of Otolaryngology--Head and Neck Surgery. Vestibular testing information.
\end{enumerate}

\clearpage
\section*{Understanding Results and Follow-up}
The laboratory report should identify the specimen, method, units, reference interval or decision limit, and whether the result is positive, negative, abnormal, or inconclusive. Reference intervals vary with age, sex, pregnancy, specimen type, assay, and laboratory; a result outside a range is not automatically a diagnosis, and a result within a range does not always exclude disease. Discuss unexpected results with the ordering clinician, who can decide whether repeat or confirmatory testing, treatment, or referral is appropriate. Do not change medicines or delay urgent care based on an isolated result.


\section*{Regulatory Status and Authorization Date}
This chapter describes a test category, not a specific commercial product. FDA, CDSCO, EU IVDR, and MHRA approval or registration dates therefore cannot be assigned without the assay manufacturer, product name, instrument, and intended use. For laboratory-developed tests, an individual agency approval date may not exist. See the Preface for the interpretation of regulatory status.

\clearpage
